\documentclass[a4paper,10pt]{article}
\usepackage[spanish]{babel}
\usepackage{
			ucs
			,amsfonts
			,amssymb
			,acronym
			,verbatim
			,rotating
			,graphicx
			,enumerate
			,fancyhdr
            ,mathrsfs
            ,array
			}

\usepackage{tikz}
\usetikzlibrary{svg.path}

\pagestyle{fancy}
\addtolength{\textheight}{1.0in}
\addtolength{\textwidth}{0.5in}
\addtolength{\voffset}{-1.0in}

\fancyhf{}
\fancyhead[L]{{\bfseries Programaci\'on I Curso 2012. Pr\'actica 10 \hfill Archivos}}
\fancyfoot[LE,RO]{\thepage}

\fancypagestyle{firstpage}{
\fancyhf{}

\fancyhead[L]{\includegraphics[width=1.75cm]{ucalp.jpg}  \\
\small{UCALP. Facultad de Ciencias Exactas e Ingenier\'ia} \\
{\bfseries Programaci\'on I Curso 2012. Pr\'actica 10 \hfill Archivos}}
\fancyfoot[LE,RO]{\thepage}
}

\thispagestyle{firstpage}

\begin{document}

~\\ \\ \\ \\

\begin{enumerate}
%1
\item Generar un archivo de texto con palabras cargadas por teclado y finalizar su llenado cuando se introduzca la palabra ``fin123''.

%2
\item Generar un archivo con un rango de enteros positivos que se ingresan por teclado. Recorrer el archivo y mostrarlos en pantalla en forma de listado.

%3
\item Generar 2 archivos: 1 con n\'umeros pares y el otro con n\'umeros impares. Luego generar un tercer archivo que contenga los n�meros de ambos archivos ordenados de menor a mayor.

%4
\item Generar un archivo de agenda que contenga los datos t\'ipicos de una persona. Realizar los procedimientos de creaci\'on, modificaci\'on y b\'usqueda por nombre y apellido.

%5
\item Realizar el ejercicio \emph{3} de la pr\'actica \emph{7.5}, leyendo los datos de entrada desde un archivo y almacenando los listados resultantes en otro/s archivo/s.

%6
\item Realizar un programa que muestre un archivo de texto, si el archivo es muy grande para entrar en la pantalla, el programa debe ``paginarlo''. \begin{small} Nota: Funcionalidad similar al programa utilitario \verb-more- (D.O.S.) o \verb-less- (Linux). \end{small}   
%7
\item Realice un par de funci\'ones/procedimientos que permitan guardar y recuperar el contenido de una lista (enlazada) en disco.

%8
\item Implemente el ejercicio \emph{2} de la pr\'actica \emph{7.5}, en donde cada sucursal entrega sus res\'umenes en un archivo separado.

%9
\item Implemente el ejercicio \emph{2} de la pr\'actica anterior (\emph{9}), pero utilizando un archivo temporal para guardar los resultados intermedios de los c\'alculos y as\'i mejorar la eficiencia. \begin{small} \emph{Nota: el archivo puede contener pares del tipo (t\'ermino, valor).} \end{small}

\end{enumerate}
\end{document}
